% !PRINTED 411
\weberpage{411}
\begin{center}
§. 132.

\gesperrt{Primitive Einheitswurzeln.}
\end{center}

Dass primitive Einheitswurzeln für jeden Grad $n$ existiren, haben wir im vorigen
Paragraphen noch nicht bewiesen. Wir müssen dies zunächst nachholen und werden dabei
auch die genaue Zahl der primitiven Einheitswurzeln feststellen.

Sei der Grad $n$ in zwei Factoren $a, b$ zerlegt, die zu einander relativ prim sind, also

\[
n = ab,
\]
und sei $\alpha$ eine $a^{\mathrm{te}}$, $\beta$ eine $b^{\mathrm{te}}$ Einheitswurzel. Dann
ist das Product

\weq{1}{r = \alpha\beta}
eine $n^{\mathrm{te}}$ Einheitswurzel. Sind $\alpha', \beta'$ zwei andere $a^{\mathrm{te}}$
und $b^{\mathrm{te}}$ Einheitswurzeln, so ist $r' = \alpha'\beta'$ auch eine
$n^{\mathrm{te}}$ Einheitswurzel, und es ist zu zeigen, dass $r'$ von $r$ verschieden ist,
wenn nicht gleichzeitig $\alpha = \alpha'$ und $\beta = \beta'$ ist.

Da nämlich $a, b$ relativ prim sind, so kann man nach \S.\,118 die ganzen Zahlen $x, y$ so
bestimmen, dass

\weq{2}{ax + by = 1}
ist; und dann folgt aus (1)

\[
\alpha = r^{by}, \quad \beta = r^{ax}.
\]
Demnach ist $\alpha$ und $\beta$ durch $r$ vollständig bestimmt, und wenn $\alpha'\beta'$
auch gleich $r$ sein soll, so muss $\alpha = \alpha'$, $\beta = \beta'$ sein.

Lässt man also in (1) $\alpha$ alle $a^{\mathrm{ten}}$, $\beta$ alle $b^{\mathrm{ten}}$
Einheitswurzeln durchlaufen, so erhält $r$ genau $ab = n$ verschiedene Werthe, und es
folgt:

I.~\gesperrt{dass in der Form $\alpha\beta$ alle $n^{\mathrm{ten}}$ Einheitswurzeln
darstellbar sind,}

und weiter:

II.~\gesperrt{dass $r$ dann und nur dann eine primitive $n^{\mathrm{te}}$ Einheitswurzel
ist, wenn $\alpha$ eine primitive $a^{\mathrm{te}}$ und $\beta$ eine primitive
$b^{\mathrm{te}}$ Einheitswurzel ist.}

Denn erstens sei $\mu$ der kleinste positive Exponent, für den $r^\mu = 1$ ist; dann ist
auch $\alpha^\mu \beta^\mu = 1$, und daraus folgt nach (2), wenn man beiderseits zur
Potenz $by$ und $ax$ erhebt,

\[
\alpha^\mu = 1, \quad \beta^\mu = 1.
\]
